\chapter{滚动窗口：指纹怎么「滚」起来}
\label{ch:gen1-rolling}

\begin{analogybox}[窗口 = 只记最近 64 个字节]
指纹 $h$ 不是「从文件头算到当前」的全体哈希（那样太慢且不好流式）。  
它只盯着\textbf{滑动的 64 字节小窗}：每前进一步，丢掉最旧的一个字节，加入最新的一个。
\end{analogybox}

\section{一步一步滚（最重要的一页）}

假设窗口宽度 $w=4$（真实是 64，这里用小数字看清过程）：

\begin{figure}[htbp]
\centering
\begin{tikzpicture}[x=0.55cm,y=0.65cm]
  \node[note] at (4,3.2) {Step A：当前窗口内的 4 个字节};
  \foreach \i/\lab in {0/A,1/B,2/C,3/D} {
    \node[byte] at (\i,2) {\lab};
  }
  \draw[decorate, decoration={brace, amplitude=4pt}] (-0.3,1.5) -- (3.3,1.5)
    node[midway,below=6pt,font=\scriptsize] {窗口};

  \node[note] at (4,0.8) {Step B：再读入新字节 E，丢掉 A};
  \node[byte, fill=gray!30] at (0,0) {A};
  \foreach \i/\lab in {1/B,2/C,3/D,4/E} {
    \node[byte] at (\i,-0.8) {\lab};
  }
  \draw[decorate, decoration={brace, amplitude=4pt}] (0.7,-1.3) -- (4.3,-1.3)
    node[midway,below=6pt,font=\scriptsize] {新窗口};

  \draw[arr] (2,1.5) -- (2,0.3);
\end{tikzpicture}
\caption{窗口右移：旧左端出，新右端进。}
\end{figure}

\section{更新公式（拆开看）}

\begin{stepbox}[Step 4 — 滚动更新三行话]
\begin{enumerate}[nosep]
  \item 把旧指纹 $h$ \textbf{左移 1 位}（二进制意义下的「滚」）；
  \item \textbf{加上}新字节通过 gear 表查到的数；
  \item \textbf{减去}被扔掉的旧字节通过 gear 表查到的数。
\end{enumerate}
\end{stepbox}

\begin{codebox}[对应代码（第一代）]
\begin{lstlisting}
// 新字节在 pos，旧字节在 pos - w
h = (h << 1) + gear[文件[pos]] - gear[文件[pos - w]];
\end{lstlisting}
\end{codebox}

\begin{figure}[htbp]
\centering
\begin{tikzpicture}[node distance=0.5cm]
  \node[gen1] (h0) {$h_{\text{旧}}$};
  \node[gen1, right=0.4cm of h0] (sh) {$\ll 1$};
  \node[gen1, right=0.4cm of sh] (add) {$+\ \mathrm{gear}[\text{新}]$};
  \node[gen1, right=0.4cm of add] (sub) {$-\ \mathrm{gear}[\text{旧}]$};
  \node[gen2, right=0.4cm of sub] (h1) {$h_{\text{新}}$};
  \draw[arr] (h0) -- (sh) -- (add) -- (sub) -- (h1);
\end{tikzpicture}
\caption{每前进 1 字节，做一次这个算术。}
\end{figure}

\section{数字例子：窗口 $w=4$ 滚一步}

\begin{table}[htbp]
\centering
\caption{玩具窗口 B-C-D-E（丢掉 A，加入 E）}
\small
\begin{tabular}{@{}lll@{}}
\toprule
\textbf{步骤} & \textbf{操作} & \textbf{说明} \\
\midrule
旧窗口 & A B C D & 当前 $h_{\text{旧}}$ 已算好 \\
读新字节 & E 进入右端 & 用 gear{[E]} \\
丢旧字节 & A 离开左端 & 用 gear{[A]} 减掉 \\
新 $h$ & $(h \ll 1) + \mathrm{gear[E]} - \mathrm{gear[A]}$ & 只改 1 次，不用重算整窗 \\
\bottomrule
\end{tabular}
\end{table}

\section{为什么要「预热」？}

刚开始扫描时，窗口里还没有凑满 $w$ 个字节。  
程序会先读够 $w$ 个，算出一个初始 $h$，这叫 \keyterm{InitWindowHash}（暖机）。

\begin{figure}[htbp]
\centering
\begin{tikzpicture}[x=0.55cm,y=0.7cm]
  \node[note, anchor=west] at (-0.2,3.3) {一个 chunk 内的字节（$w=4$，min=2 示意）};

  \foreach \i in {0,1,2,3,4,5,6,7,8,9} {
    \node[byte] at (\i,2) {\i};
  }

  \fill[CutRed!12] (-0.35,1.55) rectangle (1.65,2.55);
  \node[font=\scriptsize, CutRed!80!black] at (0.65,1.25) {禁切 min};

  \fill[Gen1Blue!18] (1.65,1.55) rectangle (5.65,2.55);
  \node[font=\scriptsize, Gen1Blue!80!black] at (3.65,1.25) {暖机：凑满窗口};

  \fill[Gen1Blue!8] (5.65,1.55) rectangle (9.65,2.55);
  \node[font=\scriptsize] at (7.65,1.25) {for 循环扫路标};

  \draw[cutline] (5.65,1.4) -- (5.65,2.7);
  \node[font=\scriptsize, align=center] at (5.65,3.05) {scan\_start\\窗口已满};

  \draw[arr, Gen1Blue] (5.65,2.85) -- (6.5,2.85)
    node[right, font=\scriptsize] {从这里 probe};
\end{tikzpicture}
\caption{for 循环不是从 chunk 头开始，而是从 min 之后、且窗口已满的 scan\_start 开始。}
\end{figure}

\begin{stepbox}[暖机 vs 主循环]
\begin{itemize}[nosep]
  \item \textbf{暖机阶段}：从 chunk 起点读 $w$ 字节，初始化 $h$（不做切点判断）；
  \item \textbf{主循环}：从 \texttt{scan\_start = pos + min} 起，每字节更新 $h$ 并检查 mask；
  \item 真实 FastCDC：$w=64$，min 通常是 256KB 量级——暖机相对整个 chunk 很短。
\end{itemize}
\end{stepbox}

\section{Probe：每走一步算一次}

工程里把「每检查一个位置能不能切」叫一次 \keyterm{probe}。  
第一代：每个字节位置 $\approx$ 1 次 probe。性能好不好，就看 probe 有多省。

\begin{figure}[htbp]
\centering
\begin{tikzpicture}[node distance=0.35cm]
  \node[byte] (p0) {pos};
  \node[byte, right=0.25cm of p0] (p1) {+1};
  \node[byte, right=0.25cm of p1] (p2) {+2};
  \node[byte, right=0.25cm of p2] (p3) {+3};
  \node[right=0.15cm of p3, font=\Large] {$\cdots$};

  \node[gen1, below=0.55cm of p0, minimum width=1.1cm] (h0) {probe};
  \node[gen1, below=0.55cm of p1, minimum width=1.1cm] (h1) {probe};
  \node[gen1, below=0.55cm of p2, minimum width=1.1cm] (h2) {probe};
  \node[gen1, below=0.55cm of p3, minimum width=1.1cm] (h3) {probe};

  \draw[arr] (p0) -- (h0);
  \draw[arr] (p1) -- (h1);
  \draw[arr] (p2) -- (h2);
  \draw[arr] (p3) -- (h3);
\end{tikzpicture}
\caption{每个候选位置：更新 $h$ → 看 mask → 命中则切，否则继续。}
\end{figure}

\section{模拟示例：$i=2$ 完整滚动}

\begin{simbox}[窗口 $w=2$，在 $i=2$ 读入 \texttt{0x30}]
\begin{tabular}{@{}ll@{}}
旧窗口 & \texttt{10, 20}（$h=18$） \\
新窗口 & \texttt{20, 30}（丢掉 \texttt{10}） \\
公式 & $h = (18 \ll 1) + \mathrm{gear[0x30]} - \mathrm{gear[0x10]}$ \\
    & $= 36 + 0 - 16 = 20$ \\
\end{tabular}
\end{simbox}
